\subsection{Alterações no Projeto}
\label{sec:alteracoes}

Em relação ao plano inicial, a representação passou de \num{60} descritores com uma rede densa para grafos dos canais ativos, porque os descritores resumem a leitura do detector inteiro; os modelos finais, porém, ainda leem \num{33} escalares do evento construídos à mão. Em vez de uma só arquitetura, foram comparados \num{11} codificadores de capacidade equiparada, para que o número de parâmetros não decidisse o resultado. \textit{Boosters} ajustados a atributos agregados entraram como referência depois que, em julho de 2026, um modelo simples de \textit{gradient boosting} superou a GNN da época. A avaliação mudou de uma partição de \qty{80}{\percent} e \qty{20}{\percent}, com validação cruzada em \num{10} partes, para uma divisão estratificada fixa de \num{7490}, \num{937} e \num{937} eventos, comum a todos os modelos; a comparação final não teve validação cruzada. O APEX foi comparado em energia igual e também em luz detectada igual, porque sua simulação tem rendimento luminoso médio de \qty{180}{\pe\per\MeV}, \num{3.4} vezes os \num{53.0} do teste.

A extensão a outras energias e geometrias, prevista no plano, não foi feita com os modelos finais; restou a varredura de eficiência e de energia do baricentro (Seção~\ref{sec:limites}). Na comparação final, os codificadores foram treinados sem a busca de hiperparâmetros do plano (Seção~\ref{sec:modelos}), aplicada só aos \textit{boosters} (Seção~\ref{sec:protocolo_comparacao}). Das métricas previstas por eixo, ficaram o RMSE, a média dos três MAE e o $R^2$ só em $z$ (Seção~\ref{sec:onde_erro}); a ordenação e a perda do treinamento usam o erro 3-D médio (Seção~\ref{sec:tarefa}). A comparação com o APEX, proposta por coordenada, ficou restrita a $z$ e à largura combinada, porque a correspondência das larguras transversais com $x$ e $y$ não está estabelecida (Seção~\ref{sec:protocolo_comparacao}). O uso em aquisição em tempo real, esperado no plano, não foi avaliado: o tempo de inferência não foi medido, e os codificadores, com capacidade equiparada a \num{1843395} parâmetros, são maiores que a rede da etapa anterior, de \num{8195}. Essa rede é citada, não reavaliada, porque usou outra partição e a geometria anterior às correções e porque o modelo treinado não está disponível.
